\documentclass{dlproblemset}
\usepackage{tikz}
\tikzset{every picture/.style={line width=0.75pt}}

\title{Conexões Residuais}
\subtitle{Lista de Exercícios}

\begin{document}

\begin{enumerate}[I)]

   \item Explique porque os dois ramos de um bloco residual não são correlacionados. Mostre que a variância da soma de duas variáveis aleatórias não correlacionadas é igual a soma das variâncias.

   \item Por que precisamos usar \textit{batch norm} quando usamos \textit{skip connections?}

   \item Quantos parâmetros treináveis uma camada de \textit{batch norm} possui?

   \item O grafo computacional do \textit{batch norm} é dado ela Figura \ref{fig:batch-norm}. Nessa figura, os circulos marcam cada uma das funções que são executadas durante o processo de um forward que aplica \textit{batch norm}. Tendo em vista esse contexto, faça o que é solicitado nas questões abaixo:

   \begin{figure}[H]
       \centering
       \includegraphics[width=1\linewidth]{problems/02-cnn/images/grafo_computacional_batch_norm.png}
       \caption{Grafo Computacional Batch Norm}
       \label{fig:batch-norm}
   \end{figure}

As equações abaixo detalham o grafo computacional (o deslocamento, que a figura escreve como $\delta$, aparece aqui como $\beta$, seguindo a notação usada em aula):

   \begin{align*}
       f_1 &= \mathbb{E}[z_i] &
       f_{2i} &= z_i - f_1 \\
       f_{3i} &= f_{2i}^2 &
       f_4 &= \mathbb{E}[f_{3i}] \\
       f_5 &= \sqrt{f_4+\epsilon} &
       f_6 &= \frac{1}{f_5} \\
       f_{7i} &= f_{2i}\, f_6 &
       z_i' &= f_{7i}\,\gamma + \beta
    \end{align*}
   \begin{enumerate}[a)]
       \item Escreva código em Python para implementar o \textit{forward pass}.
       \item Descubra as derivadas de cada nodo do grafo.
       \item Crie uma expressão que computa $\frac{\partial z'_i}{\partial z_i}$
       \item Implemente o \textit{backward pass} em Python.
   \end{enumerate}

    \item Calcule a saída do caminho $\mathbf{F}(\mathbf{X})$ e a saída final de um bloco residual onde a entrada e a saída possuem as mesmas dimensões. Lembre que:
    \begin{equation}
        \mathbf{H}(\mathbf{X}) = \mathbf{F}(\mathbf{X}) + \mathbf{X}
    \end{equation}
    Onde $\mathbf{F}(\mathbf{X})$ representa uma camada convolucional utilizando o Kernel $\mathbf{\Theta}_1$. Assuma padding = 1; stride = 1; bias = 0.

    $$\mathbf{X} = \begin{bmatrix}
    1 & 2 & 3 \\
    4 & 5 & 6 \\
    7 & 8 & 9
    \end{bmatrix}
    \qquad
    \mathbf{\Theta}_1 = \begin{bmatrix}
    0 & 0 & 0 \\
    0 & -1 & 0 \\
    0 & 0 & 0
    \end{bmatrix}
    $$

    \item Sobre a questão anterior, que aconteceria com a saída $\mathbf{H}(\mathbf{X})$ se o kernel $\mathbf{\Theta}_1$ fosse treinado para ter todos os seus 9 pesos iguais a $0$? O que isso significa para o treinamento de redes muito profundas?

    \item Em redes neurais muito profundas, os gradientes podem "desaparecer" (ficando próximos de zero), impedindo que as primeiras camadas da rede aprendam. Qual é a "chave" para resolver o problema do vanishing gradient?

\end{enumerate}

\newpage
\subsection*{Gabarito}
\color{blue}

\begin{enumerate}[I)]

   \item A rigor os dois ramos \textbf{não} são descorrelacionados em geral: uma transformação linear de $X$ é perfeitamente correlacionada com $X$, e a ReLU não elimina correlação. O que justifica a hipótese é o \textbf{momento da inicialização}: com pesos aleatórios de média zero e independentes da entrada, a saída do ramo processado não tem alinhamento sistemático com a entrada, e a covariância entre os dois ramos é aproximadamente nula. Durante o treinamento os pesos se ajustam à entrada, a covariância deixa de ser nula, e a hipótese vale apenas como motivação para a variância crescente.

   Sob essa hipótese, escrevendo $\mathrm{Var}[X] = \mathbb{E}[(X - \mathbb{E}[X])^2]$:
   \begin{align*}
       \mathrm{Var}[X+Y] &= \mathbb{E}\bigl[\bigl((X+Y) - \mathbb{E}[X+Y]\bigr)^2\bigr] \\
                          &= \mathbb{E}\bigl[(X - \mathbb{E}[X])^2\bigr] + 2\,\mathbb{E}\bigl[(X - \mathbb{E}[X])(Y - \mathbb{E}[Y])\bigr] + \mathbb{E}\bigl[(Y - \mathbb{E}[Y])^2\bigr] \\
                          &= \mathrm{Var}[X] + 2\,\mathrm{Cov}(X,Y) + \mathrm{Var}[Y].
   \end{align*}
   Com $\mathrm{Cov}(X,Y) = 0$ o termo cruzado desaparece e $\mathrm{Var}[X+Y] = \mathrm{Var}[X] + \mathrm{Var}[Y]$. Como os dois ramos têm variância $\sigma_\text{in}^2$, a saída do bloco tem $\sigma_\text{out}^2 = 2\,\sigma_\text{in}^2$, ou seja a variância dobra a cada bloco residual.

   \item Para não causar explosão das ativações nem dos gradientes devido ao aumento da variância.

   \item Dois: o fator de escala $\gamma$ e o deslocamento $\beta$. A contagem, porém, é \textbf{por unidade}: em uma camada densa com $H$ unidades ocultas são $2H$ parâmetros treináveis, e em uma camada convolucional existe um par $(\gamma, \beta)$ por \textbf{canal}, ou seja $2C$ parâmetros. A média $\mu$ e o desvio $\sigma$ também são um por unidade ou canal, mas não são treinados por \textit{backprop}: são estimados a partir do \textit{mini batch} e mantidos como médias móveis para uso em inferência.

   \item Derivadas locais de cada nodo do grafo, com $N$ o tamanho do \textit{mini batch}:
   \begin{align*}
       \frac{\partial f_1}{\partial z_i} &= \frac{1}{N} &
       \frac{\partial f_{2i}}{\partial z_i} &= 1, \quad \frac{\partial f_{2i}}{\partial f_1} = -1 \\
       \frac{\partial f_{3i}}{\partial f_{2i}} &= 2 f_{2i} &
       \frac{\partial f_4}{\partial f_{3i}} &= \frac{1}{N} \\
       \frac{\partial f_5}{\partial f_4} &= \frac{1}{2\sqrt{f_4+\epsilon}} &
       \frac{\partial f_6}{\partial f_5} &= -\frac{1}{f_5^2} \\
       \frac{\partial f_{7i}}{\partial f_{2i}} &= f_6, \quad \frac{\partial f_{7i}}{\partial f_6} = f_{2i} &
       \frac{\partial z_i'}{\partial f_{7i}} &= \gamma, \quad \frac{\partial z_i'}{\partial \gamma} = f_{7i}, \quad \frac{\partial z_i'}{\partial \beta} = 1
   \end{align*}

   Note que $z_i$ influencia $z_i'$ por três caminhos: diretamente por $f_{2i}$, pela média $f_1$ e pela variância $f_4$. Somando os três, e escrevendo $\sigma = \sqrt{f_4+\epsilon}$ e $\hat{z}_i = f_{2i}/\sigma$ para a ativação normalizada, obtemos
   \begin{equation*}
       \frac{\partial z_i'}{\partial z_j} = \frac{\gamma}{\sigma}\left(\delta_{ij} - \frac{1}{N} - \frac{\hat{z}_i\,\hat{z}_j}{N}\right),
   \end{equation*}
   em que $\delta_{ij}$ é 1 quando $i = j$ e 0 caso contrário. Em particular, o caso pedido no item (c) é a diagonal:
   \begin{equation*}
       \frac{\partial z_i'}{\partial z_i} = \frac{\gamma}{\sigma}\left(1 - \frac{1}{N} - \frac{\hat{z}_i^{\,2}}{N}\right).
   \end{equation*}
   O ponto pedagógico é que a derivada de uma ativação depende de \textbf{todo} o \textit{mini batch}, e não só da própria ativação, porque $\mu$ e $\sigma$ são estatísticas do lote. Para uma implementação de referência do \textit{forward} e do \textit{backward} em Python, itens (a) e (d), ver \href{https://kratzert.github.io/2016/02/12/understanding-the-gradient-flow-through-the-batch-normalization-layer.html}{este texto}.

        \item Cálculo de $\mathbf{F}(\mathbf{X})$ (Convolução de $\mathbf{X}$ com $\mathbf{\Theta}_1$):

O kernel $\mathbf{\Theta}_1$ é um filtro simples que multiplica o pixel central por $-1$ e todos os vizinhos por $0$. Ao passarmos $\mathbf{\Theta}_1$ sobre a matriz de entrada com padding, o resultado de cada operação é o valor do pixel central da entrada multiplicado por $-1$.

Portanto, a matriz $\mathbf{F}(\mathbf{X})$ é a negação da matriz de entrada $\mathbf{X}$:
$$
\mathbf{F}(\mathbf{X}) = \begin{bmatrix}
-1 & -2 & -3 \\
-4 & -5 & -6 \\
-7 & -8 & -9
\end{bmatrix}
$$

Cálculo $\mathbf{H}(\mathbf{X})$:
A saída final do bloco residual é a soma do resultado $\mathbf{F}(\mathbf{X})$ com a entrada original $\mathbf{X}$.
$$
\mathbf{H}(\mathbf{X}) = \mathbf{F}(\mathbf{X}) + \mathbf{X}
$$
$$
\mathbf{H}(\mathbf{X}) = \begin{bmatrix}
-1 & -2 & -3 \\
-4 & -5 & -6 \\
-7 & -8 & -9
\end{bmatrix}
+
\begin{bmatrix}
1 & 2 & 3 \\
4 & 5 & 6 \\
7 & 8 & 9
\end{bmatrix}
$$
\vspace{0.5cm}
$$
\mathbf{H}(\mathbf{X}) = \begin{bmatrix}
(-1+1) & (-2+2) & (-3+3) \\
(-4+4) & (-5+5) & (-6+6) \\
(-7+7) & (-8+8) & (-9+9)
\end{bmatrix}
=
\begin{bmatrix}
0 & 0 & 0 \\
0 & 0 & 0 \\
0 & 0 & 0
\end{bmatrix}
$$

Resultado Final:
A saída do bloco residual é uma matriz nula:
$$
\mathbf{H}(\mathbf{X}) = \begin{bmatrix}
0 & 0 & 0 \\
0 & 0 & 0 \\
0 & 0 & 0
\end{bmatrix}
$$
    \item Se o kernel $\mathbf{\Theta}_1$ fosse treinado para ser uma matriz de zeros, e assumindo que o bias $B_1$ também seja zero, o resultado da convolução $\mathbf{F}(\mathbf{X})$ seria uma matriz inteiramente nula. A Saída final do bloco, $\mathbf{H}(\mathbf{X})$, seria calculada como:
    $$ \mathbf{H}(\mathbf{X}) = \mathbf{F}(\mathbf{X}) + \mathbf{X} = \begin{bmatrix} 0 & 0 & 0 \\ 0 & 0 & 0 \\ 0 & 0 & 0 \end{bmatrix} + \mathbf{X} = \mathbf{X}$$

    O bloco se tornaria, efetivamente, uma transformação identidade, simplesmente passando a entrada $\mathbf{X}$ adiante sem qualquer alteração.
    % Isso é a principal vantagem das (ResNets), resolvendo o problema da degradação (que ocorre quando adicionar mais camadas a uma rede profunda começa a piorar o desempenho). Com as conexões residuais, se uma camada adicional não for útil (ou for prejudicial ao treinamento), o otimizador pode simplesmente aprender a "ignora-la" (fazendo seu caminho $F(x)$ tender a zero), e o bloco reverte para uma identidade. Isso garante que ao adicionar mais camadas, na pior das hipóteses, não irá prejudicar o desempenho da rede.

    \item A "chave" são as conexões residuais (\textit{skip connections}), introduzidas pelas ResNets. Elas criam um "atalho" que soma a entrada ($\mathbf{X}$) diretamente à saída de um bloco de camadas ($\mathbf{F}(\mathbf{X})$). Garantindo que o gradiente sempre tenha um caminho direto (o atalho $+\mathbf{X}$) para fluir, impedindo que ele se multiplique até desaparecer.

\end{enumerate}
\end{document}
